% BHU PYQ -- Professional Exam Template
\documentclass[11pt,a4paper]{article}

\usepackage[a4paper,top=2.5cm,bottom=2.5cm,left=2.8cm,right=2.8cm,headheight=14pt]{geometry}
\usepackage{amsmath,amssymb}
\usepackage[T1]{fontenc}
\usepackage[utf8]{inputenc}
\usepackage{lmodern}
\usepackage{microtype}
\usepackage{fancyhdr}
\usepackage{enumitem}
\usepackage{hyperref}
\usepackage{lastpage}
\usepackage{parskip}

\hypersetup{colorlinks=true,urlcolor=black,linkcolor=black,
  pdftitle={B.Sc. (Hons.) Botany Sem-V BOB-502 -- BHU PYQ}}

\pagestyle{fancy}
\fancyhf{}
\renewcommand{\headrulewidth}{0.4pt}
\renewcommand{\footrulewidth}{0.4pt}
\fancyhead[L]{\small   \textit{Banaras Hindu University}}
\fancyhead[R]{\small   \textit{Botany Paper: BOB-502}}
\fancyfoot[C]{\small Page  \thepage\ of \pageref{LastPage}}


\newlist{parts}{enumerate}{2}
\setlist[parts,1]{label=(\alph*),leftmargin=*,itemsep=4pt,topsep=3pt}
\setlist[parts,2]{label=(\roman*),leftmargin=*,itemsep=2pt,topsep=2pt}

\newcommand{\pts}[1]{\hfill   \textit{[#1]}}


\begin{document}


\begin{center}
  {\large\bfseries BANARAS HINDU UNIVERSITY}\\[2pt]
  {\normalsize B.Sc. (Hons.) Semester V Examination, 2022-23}\\[6pt]
  \rule{0.8\linewidth}{0.5pt}\\[6pt]
  {\large\bfseries Botany}\\[4pt]
  {\large\bfseries Paper: BOB-502 --- Comparative Studies of Phanerogams}\\[4pt]
  {\normalsize Time: Three Hours \hfill Maximum Marks: 70}
\end{center}

\rule{\linewidth}{0.4pt}

\smallskip



\noindent   \textit{Instructions: Answer five questions in all, selecting two each from Sections A and B and Question 9 (Section-C) which is compulsory. The figures in the right-hand margin indicate marks.}

\bigskip


\begin{center}
   \textbf{SECTION A}
\end{center}
\medskip

\medskip


\noindent   \textbf{Question 1.} Give a detailed account of reproduction in   \textit{Gnetum}.\pts{15}

\medskip


\noindent   \textbf{Question 2.} Write short notes on any two of the following:\pts{$2  \times 7.5 = 15$}

\begin{parts}
  \item Distribution of gymnosperms in India
  \item   \textit{Pentoxylon}
  \item Male and female cones of   \textit{Biota}
  \item Embryogeny in   \textit{Ginkgo}
\end{parts}

\medskip


\noindent   \textbf{Question 3.} Give a brief account of morphology and reproduction in   \textit{Zamia}.\pts{15}

\medskip


\noindent   \textbf{Question 4.} Summarize phylogenetic trends in gymnosperms.\pts{15}

\bigskip

\begin{center}
   \textbf{SECTION B}
\end{center}
\medskip

\medskip


\noindent   \textbf{Question 5.} Discuss Hutchinson's classification, its merits and demerits.\pts{15}

\medskip


\noindent   \textbf{Question 6.} Discuss briefly any two of the following:\pts{$2  \times 7.5 = 15$}

\begin{parts}
  \item Numerical taxonomy
  \item Embryo culture
  \item Parthenogenesis
  \item Cleavage polyembryony
\end{parts}

\medskip


\noindent   \textbf{Question 7.} Write salient features of the family Euphorbiaceae and discuss its economic importance.\pts{15}

\medskip


\noindent   \textbf{Question 8.} Differentiate between any three of the following:\pts{$3  \times 5 = 15$}

\begin{parts}
  \item Diplospory and Apospory
  \item Ray and Disc florets
  \item Scrophulariaceae and Verbenaceae
  \item Zingiberaceae and Liliaceae
\end{parts}

\bigskip

\begin{center}
   \textbf{SECTION C}
\end{center}
\medskip

\medskip


\noindent   \textbf{Question 9.} Answer each of the following within 50 words:\pts{$10  \times 1 = 10$}

\begin{parts}
  \item Name any monocot plant whose leaves do not show parallel venation.
  \item Write about hesperidium.
  \item Name the scientists who first reported the development of plant from anther-culture.
  \item Mention any two resemblances of   \textit{Gnetum} leaf with that of angiosperms.
  \item Who reconstructed the genus   \textit{Williamsonia sewardiana}?
  \item What do you mean by syndetocheilic stomata?
  \item Mention peculiar features of   \textit{Cuscuta} embryo.
  \item Give botanical names of any two medicinal plants associated with the Combretaceae family.
  \item Write about Cyathium inflorescence.
  \item What is Ochrea?
\end{parts}

\vfill
\rule{\linewidth}{0.4pt}

\begin{center}\small
\href{https://bhu-exam.vercel.app/}{Click Here}\\[4pt]\href{https://me-aryan.vercel.app/}{Click Here}\\[4pt]\href{https://quashan-me.vercel.app/}{Click Here}
\end{center}

\end{document}
